\chapter{Аудит селектора \texorpdfstring{$R_2$}{R2}-компоненты
Pati--Salam \texorpdfstring{$\Sigma$}{Sigma}}
\label{ch:v10-h15-r2-pati-salam-sigma-selector-audit}

Полный $\Sigma=(\mathbf2_R,\mathbf2_L,\mathbf{15}_4)$ содержит восемь
секторов Стандартной модели. В порядке предыдущего гейта оператор
$Q=6Y$ имеет спектр
\begin{equation}
 (3,-3,7,1,-1,-7,3,-3).
\end{equation}
Только $R_2$ и его сопряжение имеют $Q^2=49$; остальные значения квадрата
равны $1$ или $9$. Поэтому точный проектор получается спектральной
интерполяцией:
\begin{equation}
 \boxed{P_{R_2}=\frac{(Q^2-I)(Q^2-9I)}{1920}}.
\end{equation}
Он равен $\operatorname{diag}(0,0,1,0,0,1,0,0)$, имеет ранг $2$ и
нуллитет $6$.

Степень минимальна. Аффинный полином $a+bQ^2$ должен обращаться в нуль при
$Q^2=1$ и $9$, что принуждает $a=b=0$ и не позволяет получить единицу при
$49$. Квадратичная интерполяция единственна и даёт коэффициенты
$(9,-10,1)/1920$ в базисе $(1,Q^2,Q^4)$.

Проектор коммутирует с гиперзарядом и Real-сопряжением, обменивающим
$R_2\leftrightarrow\overline R_2$. Но его коммутатор с $SU(2)_R$-flip
имеет ранг $4$: полином каноничен лишь после выбора SM-вложения и нарушения
Pati--Salam-симметрии.

Идеальный положительный mass-split имел бы гессиан
\begin{equation}
 H_{\rm sel}=2(I-P_{R_2}),
\end{equation}
с рангом $6$ и ядром ровно на $R_2\oplus\overline R_2$. Однако текущий
унаследованный гессиан такого расщепления имеет ранг $0$. Проектор выведен
алгебраически из зарядов, но динамический коэффициент и физический масштаб
не получены.

Аудит одиннадцати маршрутов по шести критериям имеет ранг $6$ и проходы
$0/11$. Спектральный гиперзарядовый полином получает $5/6$, проваливая
только динамическое расщепление. Target-loaded complement penalty также
получает $5/6$, но не наследуется.

\begin{theorem}[Минимальный гиперзарядовый проектор]
После ограничения Pati--Salam $\Sigma$ к группе Стандартной модели
$R_2\oplus\overline R_2$ является точным спектральным подпространством
$Q^2=49$. Его единственный минимальный проектор есть квадратичный полином
по $Q^2$, приведённый выше; он SM- и Real-совместим.
\end{theorem}

\begin{nogocriterion}[Алгебраический проектор не равен mass-split]
Спектральное функциональное исчисление определяет подпространство, но не
создаёт член потенциала $\Sigma^\dagger(I-P_{R_2})\Sigma$. Пока его знак,
коэффициент и масштаб не выведены общим родителем, companion-компоненты
нельзя объявлять тяжёлыми или отсутствующими.
\end{nogocriterion}

\begin{protocol}\begin{enumerate}
\item distinct значения $Q^2$: \textbf{$1,9,49$};
\item минимальная степень по $Q^2$: \textbf{$2$};
\item ранг/нуллитет $P_{R_2}$: \textbf{$2/6$};
\item Real-совместимость: \textbf{да};
\item полный $SU(2)_R$-дефект: \textbf{rank $4$};
\item идеальный mass-Hessian rank/nullity: \textbf{$6/2$};
\item inherited mass-splitting rank: \textbf{$0$};
\item прошедших кандидатов: \textbf{$0/11$};
\item physical origin: \textbf{$1/3$};
\item следующий гейт: \textbf{родитель mass-splitting гиперзарядового проектора}.
\end{enumerate}\end{protocol}

Машинный сертификат сохранён в
\path{s2t_v10_particle_wrinkle_dislocation_callias_equal_charge_twist_m4_h15_r2_pati_salam_sigma_component_selector_candidate_audit_gate_results.json}.